% Owner: cyj. Evidence: experiments/cyj and problem/cyj audits dated 2026-09-24.
\section{问题二：分层标度关系与可识别性}

\subsection{附件角色与统一坐标的边界}

我们先按来源和生成方式固定数据角色。B1 的 1176 条 Pythia 轨迹记录用于同源 $N$--$D$ 基线；B2 为半合成 Cerebras 轨迹，B3 为插值轨迹，只用于形状检查。B4 的 57 个跨族收敛点和 B5 的 44 个文献整理点尚不能证明与 B1 使用相同 tokenizer、验证语料、对数底和聚合定义，因此不计算统一的外部 RMSE。B6 的 360 点全部包含在 B7 的 450 个 $N,D,Q$ 坐标中，两表不互作独立检验。B8 含 calibrated 与 extrapolated 标签，质量方向和损失水平与 B7 冲突，单独审计。B9 是大模型元数据，B10 是 estimated Loss，仅作为大规模外推的压力参考，不能当真实独立测试。

全文取 $N$ 的单位为十亿参数，$D$ 的单位为十亿 token。$Q_{\rm score}$ 为 B 侧原生半合成质量坐标，问题一的 $Q_A$ 是 A 侧质量代理；附件没有成对标定样本。问题一配比 $\boldsymbol p$ 对应 17 域单纯形，其可核验输出是 1M 条件下 13 个目标域的相对 Loss 对比，并非 B1 或 B7 的总体 \texttt{val\_loss}。

\subsection{B1 同源 $N$--$D$ 经验基线}

在 B1 的范围 $N\in[0.070542,11.965825]$、$D\in[0.134,299.893]$ 内，采用待检验的经验式
\begin{equation}
\widehat L_{B1}(N,D)=E+A N^{-\alpha}+B D^{-\beta}.
\end{equation}
以 8 个规模组等权、对数 Loss 残差的 Huber 目标和正参数约束估计得到
$E=1.689838$、$A=0.353969$、$B=1.240275$、$\alpha=0.339985$、$\beta=0.279892$。
全样本 RMSE 为 $1.4658\times10^{-4}$；留一规模组均值为 $1.4613\times10^{-4}$；每组后 30\% token 检查点留出的 RMSE 为 $1.1600\times10^{-4}$。去除 $E$、$N$ 或 $D$ 项均明显损害 B1 同源重构。

上述接近确定性的精度触发来源审查。B1 的 $D$ 与训练步数乘每步 $2{,}097{,}152$ token 的显示舍入相容，\texttt{ppl} 与 $\exp(\texttt{val\_loss})$ 的显示舍入也全行相容；但本附件缺少逐检查点原始评估记录和生成脚本，尚不能判断 \texttt{val\_loss} 是否经插值、平滑或共同公式构造。因此这里只主张\emph{B1 同源数据上的条件经验重构}，不以极低残差推断真实跨源预测精度。

\subsection{B7 原生质量关系及交互诊断}

B7 的原生支持为 $N\in[0.07,11.97]$、$D\in[10,600]$、$Q_{\rm score}\in[0.1,1]$。已发布的诊断基线为
\begin{equation}
\widehat L_{B7}=E_7+A_7N^{-\alpha_7}+B_7D^{-\beta_7}+G(1-Q_{\rm score}),
\end{equation}
其 $\partial\widehat L_{B7}/\partial Q_{\rm score}=-G$ 恒定。对 45 个固定 $(N,D)$ 组分别以组内十个 Q 点回归，斜率范围为 $[-0.6072,-0.1931]$，对恒定斜率的组级 RMSE 为 $0.1078$。用 $\log N$ 与 $\log D$ 描述这些斜率的回归 $R^2=0.8601$；这只描述半合成网格，不是质量的因果机制。

预先冻结五个简洁候选后，以留 N、留 D、留 Q 等级的嵌套分组比较发现，$Q\times\log N$、$Q\times\log D$ 双交互形式在 24 个外层折中的 22 折被内层选中，另 2 折选单 $Q\times\log N$。双交互与恒定 $G$ 的平均 RMSE 在 N 留级为 $0.05194$ 对 $0.05823$，D 留级为 $0.05021$ 对 $0.05709$，Q 留级为 $0.05006$ 对 $0.05762$。由于候选形式在此前的全 B7 诊断后提出，这些分数仍属同源探索性比较；现阶段保留恒定 $G$ 为软件诊断基线，不将交互参数当已验证真实质量规律。

\subsection{冲突、外推及联合模型结论}

B8 的 calibrated 90/90 组和 extrapolated 60/60 组均表现为 Q 增大、Loss 上升，与 B7 的 45/45 组相反。B7/B8 共有 224 个精确 $N,D,Q$ 坐标而 Loss 全不相同；B8 还有 362/1704 个精确 $0.5$ 值，生成机制未给出。我们保持 B8 隔离，不反转 Q 后强行并表。B10 的 128 个 estimated Loss 点对应 B9 元数据，且参数规模全部超出 B1 的 N 范围；它们只用于压力参考，不报告独立外推误差。

现有附件没有在统一 Loss 坐标下同时观测 $N,D,Q_{\rm score},\boldsymbol p$。A 侧 $Q_A$ 到 B 侧 Q、13 个目标 Loss 到 B1/B7 Loss、配比效应跨规模变化的参数都缺少成对识别证据。因此完整 $L(N,D,Q,\boldsymbol p)$ 当前\textbf{不可识别}。后续预算决策使用分层的 B 原生诊断预测与 A 侧多目标配比敏感性；不构造未经验证的四维结构律。已有参数 bootstrap 仅覆盖固定模型的条件均值，模型形式、观测残差、跨来源以及 Loss--Benchmark 桥接的不确定性未形成总预测区间。
